\section{Compassion and empathy: what happens if we do nothing?}
\label{sec:compassion}

Begin where every successful health law begins, with a person. A lawmaker may
hear statistics from experts all day, but a patient describing a long wait for a
trial slot, or a parent describing a child with a rare cancer, creates an
emotional anchor that the data alone cannot. The honest question a patient story
poses is the one a committee cannot avoid: what happens if we do nothing?

[DRAFTING INSTRUCTIONS: in the full-narrative, ground this section in
\texttt{review/references/author\_works.bib} entry
\texttt{Kawchak2026CancerPatientsJourney} (A Cancer Patient's Journey Through a
Regulated Physical AI Oncology Trial, DOI 10.5281/zenodo.19119939) and pair the
appeal with credible evidence from \texttt{review/references/narrative\_refs.bib}
entry \texttt{novak2020patient}. Reproduce Mermaid figure
\texttt{review/mermaid/diagrams/11-patient-journey.md} as a colored TikZ figure.
Add a body-width table of the four patient archetypes (serious illness, family
caregiver, child with rare disease, person unable to access treatment) with the
column format \texttt{>\{\textbackslash raggedright\textbackslash arraybackslash\}p\{\textless width\textgreater\}}.]

\begin{narrativefig}
\node[nfone] (a) at (0,0) {Patient and family};
\node[nftwo] (b) at (3.8,0) {Verified trial access};
\node[nfhope] (c) at (7.8,0) {Protected outcome};
\draw[nfedge] (a)--(b); \draw[nfedge] (b)--(c);
\end{narrativefig}
\vspace{-0.2cm}
\figcaption{Figure 2. The compassion path (placeholder; expanded from Mermaid 11
in the full-narrative).}

The point is not to substitute feeling for evidence. It is to make the evidence
memorable by attaching it to the people it protects, which is exactly the balance
the advocacy literature recommends.
